\section*{Capítulo \ref{ch:b1:hotspots} --- Puntos calientes y vulcanismo intraplaca}

\begin{solution}{exo:b1:hotspots:1}
(a) $4205 + 5274 = \qty{9479}{m}$, casi \qty{9.5}{km}. El volcán es más pesado que el agua que desaloja: la placa se curva hacia abajo bajo su carga,
así que su base yace más abajo que el fondo lejano, y su altura real es mayor.
(b) $74000/0.1 = \num{740000}$ años con el suministro del Kilauea;
$74000/0.025 = \num{3.0e6}$ años con la velocidad media de la cadena. (c)
$0.12\times19 = \qty{2.3}{km^3}$, unas $\num{3e-5}$ del Mauna Loa: tres
milésimas de uno por ciento.
\end{solution}

\begin{solution}{exo:b1:hotspots:2}
(a) $0.92\times\pi/180 = \qty{0.0161}{rad.Myr^{-1}}$. (b) En radianes por
millón de años por kilómetros, $\omega R\sin\Delta$ está en kilómetros por
millón de años, es decir milímetros al año: $0.016057\times6371 = 102.3$, así que
$v = 102.3\sin\Delta$, lo que da \qty{51.1}{mm.yr^{-1}} a \ang{30},
\qty{88.6}{mm.yr^{-1}} a \ang{60}, \qty{102.3}{mm.yr^{-1}} a \ang{90} y
\qty{97.3}{mm.yr^{-1}} a \ang{108.0}. (c) La velocidad es proporcional a la
distancia $R\sin\Delta$ al eje, máxima sobre el \omterm{def:b1:spherical-astronomy:spherical-triangle}{círculo máximo} a \ang{90}
del polo; más allá, el punto se acerca otra vez al eje, del lado del polo
opuesto.
\end{solution}

\begin{solution}{exo:b1:hotspots:3}
(a) Deccan: $\num{e6}/\num{5e5} = \qty{2}{km}$. Atlántico Central:
$\num{2e6}/\num{7e6} = \qty{0.29}{km}$, unos \qty{290}{m}. (b) La media para
el Deccan se acerca a su sección de \qty{2.5}{km}: la provincia es una meseta gruesa y continua. (c) El área del Atlántico Central es la de toda la región, en cuatro continentes, sobre la que se encuentran sus restos: lavas, diques y filones capa hoy separados por un océano y muy erosionados. La roca se concentra en algunos lugares y falta en la mayor parte del área, así que la media no dice nada de una sola pila.
\end{solution}

\begin{solution}{exo:b1:hotspots:4}
(a) $\rho_m\alpha\,\Delta T = 3300\times\num{3.2e-5}\times250 =
\qty{26.4}{kg.m^{-3}}$. (b) La columna del abultamiento tiene la masa en exceso
$(\rho_m - \rho_w)h = 2300\times690 = \qty{1.587e6}{kg.m^{-2}}$, que la
capa ligera debe equilibrar: $H = \num{1.587e6}/26.4 = \qty{6.0e4}{m}$, unos
\qty{60}{km} de roca caliente.
\end{solution}

\begin{solution}{exo:b1:hotspots:5}
(a) En A ($\lambda = 0$, $\varphi = 0$), $\vec e = (0, 1, 0)$ y
$\vec n = (0, 0, 1)$; con $\sin\ang{270} = -1$ y $\cos\ang{270} = 0$,
$\vec t_A = -\vec e = (0, -1, 0)$. En B ($\lambda = \ang{90}$),
$\vec e = (-1, 0, 0)$, $\vec n = (0, 0, 1)$ y $\vec t_B = (1, 0, 0)$.
(b) $\vec t_A\times\vec t_B = (0, 0, 1)$. (c) Con $\vec p' = (0, 0, 1)$ y
$\vec r_A = (1, 0, 0)$, $\vec p'\times\vec r_A = (0, 1, 0)$, y
$(\vec p'\times\vec r_A)\cdot\vec t_A = -1 < 0$: el polo es
$-\vec p' = (0, 0, -1)$, el polo sur. Un movimiento hacia el oeste a lo largo del ecuador
es horario visto desde encima del polo norte, así que antihorario visto desde encima del
polo sur. (d) A \ang{90} del polo, $\omega = 50/6371 =
\qty{7.85e-3}{rad.Myr^{-1}} = \qty{0.45}{\degree\per\Myr}$.
\end{solution}

\begin{solution}{exo:b1:hotspots:6}
(a) $\sum x_i v_i = 90000 + 240000 + 342000 = \num{672000}$ y
$\sum x_i^2 = \num{7.0e7}$, así que $\omega = \qty{0.0096}{rad.Myr^{-1}}$, o
\qty{0.55}{\degree\per\Myr}. (b) Las predicciones son 28.8, 48.0 y
\qty{57.6}{mm.yr^{-1}}; los residuos $+1.2$, $0.0$ y $-0.6$
\unit{mm.yr^{-1}}. (c) Para una rotación rígida la velocidad es proporcional a la
distancia al eje, y un punto sobre el eje no se mueve: no hay
ordenada en el origen que ajustar.
\end{solution}

\begin{solution}{exo:b1:hotspots:7}
(a) $v = \qty{50e-3}{m}/\qty{3.156e7}{s} = \qty{1.58e-9}{m.s^{-1}}$, así que
$B = 2300\times\num{3e8}\times\num{1.58e-9} = \qty{1090}{kg.s^{-1}}$ y
$Q = (910/\num{3.2e-5})\times1090 = \qty{3.1e10}{W}$. (b) Un tercio del
flujo hawaiano. (c) $B$ crece con $A\,v$: la pluma bajo la placa más rápida debe rehacer el mismo abultamiento el doble de deprisa, y su flujo es el doble.
\end{solution}

\begin{solution}{exo:b1:hotspots:8}
(a) La velocidad es $350/(16.1 - 10.3) = 350/5.8 = \qty{60}{km.Myr^{-1}}$, es decir
\qty{60}{mm.yr^{-1}}. (b) La traza joven, \qty{29 \pm 5}{mm.yr^{-1}}, concuerda
con los \qty{26 \pm 5}{mm.yr^{-1}} de la placa dentro de las incertidumbres; la
parte antigua fue más del doble de rápida, $60 - 26 = 34$, casi siete incertidumbres estándar
de la velocidad de la placa. (c) Sí: N\,\ang{54}\,E es el \omterm{def:b1:spherical-astronomy:horizontal}{acimut}
\ang{54}, opuesto a \ang{235} dentro de un grado. El vulcanismo se mueve sobre la
placa en sentido opuesto a la placa sobre el \omterm{def:b1:hotspots:hotspot}{punto caliente}. (d) El estiramiento separó
las calderas más antiguas, alargando las distancias entre ellas y con ello la
velocidad aparente; la placa en sí no fue más rápida.
\end{solution}

\begin{solution}{exo:b1:hotspots:9}
(a) $\num{9.2e10}/\num{e12} = \qty{0.092}{W.m^{-2}}$, \qty{92}{mW.m^{-2}}.
(b) Entre nueve y dieciocho veces el exceso observado, que incluso podría ser nulo.
(c) $L^2/\kappa =
\num{e10}/\num{e-6} = \qty{e16}{s}$, unos \qty{320}{Myr}; la placa cruza
\qty{1000}{km} en $1000/92 = \qty{11}{Myr}$. El calor desde abajo ha tenido tiempo de
cruzar solo una pequeña parte de la placa mientras esta yace sobre el abultamiento: la mayor parte del
calor de la pluma se queda en el \omterm{def:b1:layered-earth:layers}{manto} bajo la placa y se lleva con
ella, y la superficie ve poco de él.
\end{solution}

\begin{solution}{exo:b1:hotspots:10}
(a) Diez veces el fundido, \qty{e7}{km^3}, de radio
$(3V/4\pi)^{1/3} = \qty{134}{km}$. (b) $\num{e7}/100 = \qty{e5}{km^2}$, un disco
de diámetro $2\sqrt{\num{e5}/\pi} = \qty{357}{km}$: dieciocho veces menor en área
que las trampas. La cabeza fue mayor o se fundió menos, o se extendió más y más
delgada, o las lavas fluyeron mucho más allá de ella. (c) \qty{1}{km^3} al año, diez
veces el suministro actual del Kilauea. (d) Una pluma nueva debe reunir roca caliente y ligera
suficiente para abrirse paso hacia arriba a través del \omterm{def:b1:layered-earth:layers}{manto} viscoso, y un volumen grande
asciende mucho más deprisa que uno pequeño; la cola que sigue fluye hacia arriba por el camino caliente y débil que ha abierto la cabeza, y basta un conducto estrecho.
\end{solution}

\begin{solution}{exo:b1:hotspots:11}
(a) $\qty{40}{km.Myr^{-1}}\times(81 - 47)\,\unit{Myr} = \qty{1360}{km}$.
(b) La placa se mueve a $(-80\sin\ang{20}, 80\cos\ang{20}) = (-27.4, 75.2)$
y el \omterm{def:b1:hotspots:hotspot}{punto caliente} en $(0, -40)$: respecto del \omterm{def:b1:hotspots:hotspot}{punto caliente} la placa se mueve a
$(-27.4, 115.2)$ \unit{mm.yr^{-1}}, una velocidad de \qty{118.4}{mm.yr^{-1}} hacia
$\ang{360} - \arctan(27.4/115.2) = \ang{346.6}$. (c) La traza daría una
velocidad \qty{48}{\percent} demasiado alta y una dirección \ang{6.6} demasiado al este:
un \omterm{def:b1:hotspots:hotspot}{punto caliente} que deriva registra el movimiento de la placa respecto de sí mismo, no respecto del \omterm{def:b1:layered-earth:layers}{manto} profundo.
\end{solution}

\begin{solution}{exo:b1:hotspots:12}
(a) $\qty{3e6}{J.m^{-3}.K^{-1}}\times\qty{1300}{K}\times\qty{e15}{m^3} =
\qty{3.9e24}{J}$. (b) En $\qty{3.16e13}{s}$, \qty{1.2e11}{W}: casi el calor que transporta
hoy la pluma hawaiana (1.3 veces), y \qty{0.26}{\percent} de los \qty{47}{TW} de la Tierra. (c) Una cota inferior: la cabeza también
eleva el calor en exceso de los nueve décimos de su roca que no se funde, y calienta la base de la \omterm{def:b1:layered-earth:lithosphere}{litosfera}; el calor latente se ha dejado fuera.
\end{solution}

\begin{solution}{pb:b1:hotspots:1}
\textbf{1.} En Hawái ($\varphi = \ang{19.0}$, $\lambda = \ang{-155.2}$),
con $\sin\ang{304} = -0.8290$ y $\cos\ang{304} = 0.5592$:
\[
\begin{aligned}
  \vec r_H &= (-0.8583,\ -0.3966,\ 0.3256), &
  \vec e &= (0.4195,\ -0.9078,\ 0),\\
  \vec n &= (0.2955,\ 0.1366,\ 0.9455), &
  \vec t_H &= (-0.1825,\ 0.8289,\ 0.5287).
\end{aligned}
\]

\textbf{2.} En Louisville ($\varphi = \ang{-53.6}$, $\lambda = \ang{-140.6}$),
con $\sin\ang{316} = -0.6947$ y $\cos\ang{316} = 0.7193$:
\[
\begin{aligned}
  \vec r_L &= (-0.4586,\ -0.3767,\ -0.8049), &
  \vec e &= (0.6347,\ -0.7727,\ 0),\\
  \vec n &= (-0.6220,\ -0.5109,\ 0.5934), &
  \vec t_L &= (-0.8883,\ 0.1693,\ 0.4269).
\end{aligned}
\]

\textbf{3.} $\lVert\vec t_H\rVert^2 = 0.0333 + 0.6871 + 0.2795 = 0.9999$ y
$\vec t_H\cdot\vec r_H = 0.1566 - 0.3287 + 0.1721 = 0.0000$, dentro del redondeo.

\textbf{4.} $\vec t_H\times\vec t_L = (0.2643, -0.3917, 0.7054)$, de longitud
$0.8491$. El ángulo entre las trazas es $\arcsin0.8491 = \ang{58.1}$ (su producto
escalar, $0.528$, es positivo, así que el ángulo es agudo).

\textbf{5.} $\vec p' = (0.3113, -0.4613, 0.8308)$. Entonces
$(\vec p'\times\vec r_H)\cdot\vec t_H = -0.98 < 0$: el polo es
$\vec p = (-0.3113, 0.4613, -0.8308)$.

\textbf{6.} $\arcsin(-0.8308) = \ang{-56.18}$ y
$\operatorname{atan2}(0.4613, -0.3113) = \ang{124.01}$: el polo yace en
\ang{56.2}\,S, \ang{124.0}\,E, en el océano Austral al sur de Australia.

\textbf{7.} $\vec p\cdot\vec r_H = 0.2672 - 0.1830 - 0.2705 = -0.1863$, así que
$\Delta_H = \ang{100.7}$; $\vec p\cdot\vec r_L = 0.1428 - 0.1738 + 0.6687 =
0.6377$, así que $\Delta_L = \ang{50.4}$.

\textbf{8.} El vector $\vec p\times\vec r_L$ es perpendicular a $\vec p$ y
a $\vec r_L$, como lo es $\vec t_L$ (por construcción de $\vec p$, y porque una
traza es horizontal); ambos yacen sobre la misma recta, y su longitud es
$\sin\Delta_L$. Así que $(\vec p\times\vec r_L)\cdot\vec t_L = \pm\sin\Delta_L$, y
el cálculo da $+0.770 = \sin\ang{50.4}$: la placa se mueve en Louisville
a lo largo de $+\vec t_L$, en el sentido en que corre la traza. La regla de signos, aplicada en Hawái, podría haber dado el sentido equivocado en Louisville; no lo hizo, así que ambas direcciones se reproducen exactamente, \ang{304} y \ang{316}.

\textbf{9.} $\omega = 67/(6371\sin\ang{50.38}) = 67/4908 =
\qty{0.01365}{rad.Myr^{-1}}$, es decir \qty{0.782}{\degree\per\Myr}.

\textbf{10.} $360/0.782 = \qty{460}{Myr}$.

\textbf{11.} La regla del coseno entre \ang{56.18}\,S, \ang{124.01}\,E y
\ang{63.78}\,S, \ang{112.48}\,E da \ang{9.5}. La velocidad de dos trazas,
\qty{0.78}{\degree\per\Myr}, es un \qty{15}{\percent} menor que los
\qty{0.92}{\degree\per\Myr} ajustados a trece velocidades.

\textbf{12.} Sociedad \ang{73.6}, Pitcairn \ang{77.8}, Cobb \ang{134.7}.

\textbf{13.} La velocidad $\omega R = 0.01365\times6371 = \qty{86.98}{mm.yr^{-1}}$, por
$\sin\Delta$: Sociedad \qty{83.5}{mm.yr^{-1}}, Pitcairn
\qty{85.0}{mm.yr^{-1}}, Cobb \qty{61.8}{mm.yr^{-1}}.

\textbf{14.} Pitcairn concuerda: \qty{5}{mm.yr^{-1}} menos que lo medido, dentro de
su \qty{15}{mm.yr^{-1}}. Sociedad es \qty{25.5}{mm.yr^{-1}} demasiado lenta, 2.6
incertidumbres, y Cobb \qty{18.8}{mm.yr^{-1}} demasiado rápida, más de seis:
ninguna concuerda.

\textbf{15.} El polo es perpendicular a ambos vectores de traza. Si son
casi paralelos, su producto vectorial es corto, y un error de unos grados en
cualquier sentido lo gira un ángulo grande: el polo está mal determinado.

\textbf{16.} $v = \qty{92e-3}{m}/\qty{3.156e7}{s} = \qty{2.92e-9}{m.s^{-1}}$
y $B = 2300\times\num{4.83e8}\times\num{2.92e-9} = \qty{3240}{kg.s^{-1}}$.

\textbf{17.} $Q = (910/\num{3.2e-5})\times3240 = \qty{9.2e10}{W}$,
\qty{0.20}{\percent} de \qty{47}{TW}.

\textbf{18.} Las dorsales producen $21/67000 = \qty{3.1e-4}{km^3}$ al año por
kilómetro; $0.1/\num{3.1e-4} = \qty{320}{km}$ de dorsal.

\textbf{19.} $74000/0.1 = \num{740000}$ años.

\textbf{20.} $\num{2e6}/\num{81e6} = \qty{0.025}{km^3}$ al año: el Kilauea solo
recibe hoy cuatro veces la producción media de todo el \omterm{def:b1:hotspots:hotspot}{punto caliente}, que está
cerca de su máximo.

\textbf{21.} $v_H = 86.98\sin\ang{100.73} = 86.98\times0.9825 =
\qty{85.5}{mm.yr^{-1}}$.

\textbf{22.} $92 - 85.5 = \qty{6.5}{mm.yr^{-1}}$, 2.2 incertidumbres estándar de
la velocidad medida: los dos \omterm{def:b1:hotspots:hotspot}{puntos calientes} no concuerdan del todo sobre el movimiento de una sola \omterm{def:b1:plate-kinematics:rigid-plate}{placa rígida} sobre un marco fijo.

\textbf{23.} Las edades dan \qty{9.1}{cm} y \qty{8.6}{cm} al año, promediadas sobre toda la cadena; el \qty{9.2}{cm} al año de Morgan y Morgan
se acerca al primero, y la predicción de Louisville, \qty{8.5}{cm} al año, al segundo. Las cuatro están a menos de \qty{0.7}{cm} al año unas de
otras: trazas, edades y direcciones cuentan la misma historia con un margen de unos \qty{8}{\percent}.

\textbf{24.} Con un decimal, la traza de Louisville da a la placa del Pacífico sobre el \omterm{def:b1:hotspots:hotspot}{punto caliente} de Hawái la velocidad
\textbf{\qty[text-series-to-math=true]{8.5}{cm.yr^{-1}}}.
\end{solution}
